\appendix

\section{Symbols}
\label{app:symbols}

Tables~\ref{tab:symbols} to \ref{tab:symbols3} list every symbol of this paper with its meaning
and the place of its definition. No symbol has two meanings. Where one letter occurs with different
indices or in different type, the tables say how the uses are told apart; a machine check of the
registry from which these tables are taken verifies that no symbol and no name in a data file is
defined twice (Appendix~\ref{app:checks}).

\begin{table*}
\caption{\label{tab:symbols}The symbols of this paper: geometry and the Dirac field.}
\begin{ruledtabular}
\begin{tabular}{p{0.2\textwidth}p{0.6\textwidth}p{0.14\textwidth}}
symbol & meaning & defined in \\
\hline
$x0,\dots,x7$ & the eight coordinates: $x0$ hidden space-like coordinate; $x1,x2,x3$ ordinary space; $x4$ ordinary time; $x5,x6,x7$ the deflating coordinates, \SIG{time-like}{space-like} & Sec.~\ref{sec:intro} \\
$xt$; $xj$; $a$ & a deflating coordinate, $t\in\{5,6,7\}$; a flat coordinate other than the time, $j\in\{0,1,2,3\}$; index of the four extra axes, $a\in\{0,5,6,7\}$ & Sec.~\ref{sec:intro} \\
$\gff$, $\gss$, $\gvv$; $\rt$ & the three rates of deflation (upright, without subscript; not components of the metric); the rate of $xt$ & Eq.~(\ref{eq:hyp}) \\
$\cff$, $\css$, $\cvv$; $\ct$ & the three positive constants of the metric: the scale factors at $x4=0$ & Eq.~(\ref{eq:hyp}) \\
$\gsum$, $\gsq$, $\gprod$ & the sum of the rates, the sum of their squares, the sum of their products in pairs & Sec.~\ref{sec:geometry} \\
$b_t$ & scale factor of the deflating coordinate $xt$ & Eq.~(\ref{eq:scale}) \\
$\etam$; $\gm_{\mu\nu}$, $\gmup{\mu\nu}$; $e_A{}^{\mu}$ & flat metric, metric of the deflating spacetime and its inverse, in the sign convention of the hypothesis (space-like components positive, time-like components negative; the superscript marks the spacetime, so that a component of the metric cannot be taken for a rate or a coupling); orthonormal frame & Eqs.~(\ref{eq:etam}), (\ref{eq:metric}) \\
$\Vd$ & proper volume of a unit coordinate cell of the deflating coordinates & Eq.~(\ref{eq:Vd}) \\
$\Gamma^{\rho}{}_{\mu\nu}$; $R^{\rho}{}_{\sigma\mu\nu}$, $R^{\mu}{}_{\nu}$, $R$ & Christoffel symbols; Riemann tensor [in Eq.~(\ref{eq:riemann}) its components in the orthonormal frame], Ricci tensor, curvature scalar & Sec.~\ref{sec:curvature} \\
$\cG^{\mu}{}_{\nu}$; $\cH^{\mu}{}_{\nu}$; $\Lam$, $\aGB$ & Einstein tensor; Lanczos tensor of the Gauss-Bonnet term; cosmological constant and coupling of Einstein-Gauss-Bonnet gravity & Sec.~\ref{sec:egb} \\
$\cX_{\mathrm{scale}}$, $\cX_{(s)}$ & the scaling Killing vector; the three special Killing vectors of the case of equal rates & Eqs.~(\ref{eq:scaling}), (\ref{eq:special}) \\
$K_{AB}$; $\bar x4$, $\bar xt$ & sectional curvature of the plane of the frame directions $A$ and $B$ (two subscripts; $K_{4i}$, $K_{ij}$ in Sec.~\ref{sec:gravity}); coordinates adapted to the case of equal rates & Appendix~\ref{app:curvature} \\
$P_\mu$; $\varsigma$\SIG{; $x4_{\mathrm{turn}}$}{} & conserved momenta of a classical particle; affine parameter\SIG{; turning value of the ordinary time}{} & Sec.~\ref{sec:geodesics} \\
$\gamma^{A}$, $\gamma^{\mu}$; $\Omega_\mu$, $\omega_{\mu AB}$ & Dirac matrices of the frame (printed in Appendix~\ref{app:matrices}) and in coordinate components; spin connection and its coefficients & Sec.~\ref{sec:dirac} \\
%%<1644>
$\Cm$; $\chi$, $P^{L,R}$ & the matrix that turns each Dirac matrix into minus its transpose; chirality matrix and the two chiral projectors & Secs.~\ref{sec:dirac}, \ref{sec:clock} \\
%%</1644>
%%<1671>
$\chi$, $P^{L,R}$; $\Cm$ & chirality matrix and the two chiral projectors; the matrix that turns each Dirac matrix into minus its transpose & Secs.~\ref{sec:clock}, \ref{sec:gap} \\
%%</1671>
$\psi_{\tg}$; $\hpsi$ & the 16-component field on the deflating spacetime; the weighted field $\sqrt{\Vd}\,\psi$ & Eqs.~(\ref{eq:dirac}), (\ref{eq:weighted}) \\
%%<1644>
$K$, $K_\perp$; $\gamma^{d}$; $W_{1644}$ & generator of the evolution of a mode in the time, and its part without the scale factors; frame matrix along the wave vector in the deflating coordinates; matrix of the conserved charge & Eq.~(\ref{eq:K}) \\
$\beta$, $\alpha^{j}$; $\bar\psi$; $\tilde\psi$ & the Hermitian matrices $-i\gamma^{4}$ and $-\gamma^{4}\gamma^{j}$ ($j=0,1,2,3$); the adjoint $\psi^\dagger\beta$ of the quantized sector; the covariant adjoint $\psi^\dagger\Cm$ & Secs.~\ref{sec:lagrangian}, \ref{sec:forms} \\
$k_j$, $k_t$; $\Eperp$; $\kd$; $\ud$ & wavenumbers along the flat and the deflating coordinates; $(m^{2}+\sum_jk_j^{2})^{1/2}$; physical wavenumber along the deflating coordinates; $\kd/\rt$ & Eqs.~(\ref{eq:kd}), (\ref{eq:ud}) \\
$w$; $\Swkb$ & one of the eight constant spinors from which the exact solutions are built; exponent of the extinction of a bounded solution & Eqs.~(\ref{eq:exact}), (\ref{eq:Swkb}) \\
%%</1644>
%%<1671>
$\hh$, $\hh_\perp$; $\beta$, $\alpha^{A}$, $\alpha^{d}$; $\bar\psi$ & Hamiltonian of a mode, and its part without the scale factors; $\gamma^{4}$, $\gamma^{4}\gamma^{A}$, and the matrix $\alpha$ along the wave vector in the deflating coordinates; the adjoint $\psi^\dagger\beta$ & Eqs.~(\ref{eq:h}), (\ref{eq:LVS}) \\
$p_j$, $k_t$; $\Eperp$; $\kd$; $\ud$ & wavenumbers along the flat and the deflating coordinates; $(m^{2}+\sum_jp_j^{2})^{1/2}$; physical wavenumber along the deflating coordinates; $\kd/\rt$ & Eqs.~(\ref{eq:kd}), (\ref{eq:ud}) \\
$w$; $Z_+$, $Z_-$ & one of the eight constant spinors from which the exact solutions are built; two cylinder functions whose orders differ by one & Eq.~(\ref{eq:exact}) \\
$\ncr$; $\Tdefl$ & number of quanta that the deflation creates in a mode, per internal state; temperature of their spectrum & Eqs.~(\ref{eq:ncr}), (\ref{eq:Tdefl}) \\
%%</1671>
\end{tabular}
\end{ruledtabular}
\end{table*}

\begin{table*}
\caption{\label{tab:symbols4}The symbols of this paper, continued: mass term and self-interaction,
the Kohn-Sham and CJT calculations, and gravity.}
\begin{ruledtabular}
\begin{tabular}{p{0.2\textwidth}p{0.6\textwidth}p{0.14\textwidth}}
symbol & meaning & defined in \\
\hline
$\mathcal L_{\tg}$; $\mathcal L_{\mathrm{kin}}$, $\mathcal L_{\mathrm{mass}}$, $\Lself$ & Lagrangian of the fermion field, and its three parts: kinetic term, mass term, self-interaction & Eq.~(\ref{eq:Ltot}) \\
$g_r$; $g_0,\dots,g_8$ & the nine couplings of the self-interaction, one for each rank (one digit as subscript: not a rate $\gff$, not a component $\gm_{\mu\nu}$ of the metric, not the coupling $\gVS$) & Eq.~(\ref{eq:Lself}) \\
$K$, $J$; $r$, $r'$, $l$ & strictly increasing lists of indices out of $0,\dots,7$; their ranks (numbers of indices), and a summation index\SIG{ (the matrix $K$ of Sec.~\ref{sec:modes} is written without index or with the index $\perp$)}{} & Eq.~(\ref{eq:Theta}) \\
$\Theta^{K}$, $\Theta_K$; $\zeta_r$; $\sgnK$ & the 256 matrices of the self-interaction, $\Theta^{K}=\zeta_r\gamma^{A_1}\cdots\gamma^{A_r}$, and their inverses $\Theta_K=\sgnK\,\Theta^{K}$; the phase $1$ or $i$; the sign $\pm1$ & Eq.~(\ref{eq:Theta}) \\
$I_r$ & the nine quartic invariants & Eq.~(\ref{eq:Lself}) \\
$\sfF$, $\sfF_{rr'}$ & the Fierz matrix (sans-serif; the amplitudes $F_S$, $F_P$, $F_0,\dots$ are in italic type) & Eq.~(\ref{eq:fierzF}) \\
$\geff0$, $\geff1$, $\geff7$, $\geff8$ & the four independent couplings: scalar, vector, axial vector, pseudoscalar & Eq.~(\ref{eq:geff}) \\
$\sfn$; $n$, $\nS$ & density matrix of matter; its number density $\langle\psi^\dagger\psi\rangle$ and scalar density $\langle\bar\psi\psi\rangle$ & Sec.~\ref{sec:hf} \\
$\GS^{\mathrm{HF}}$, $\GV^{\mathrm{HF}}$\SIG{; $\GS^{\mathrm c}$, $\GV^{\mathrm c}$}{} & the values that the Hartree-Fock approximation gives to the mean-field couplings\SIG{; the two further couplings of the covariant form}{} & Eq.~(\ref{eq:GHF}) \\
$\Gpair$ & contact coupling of a pair of two particles at rest & Eq.~(\ref{eq:Gpair}) \\
%%<1644>
$\gamma_{44q}^{A}$; $I_r^{\mathrm{cov}}$ & the matrices of the quantized-sector form of the self-interaction; the invariants of the covariant form & Sec.~\ref{sec:forms} \\
$n_\pm$, $n_{S\pm}$; $n_W$, $n_{SW}$; $\mathrm{pol}_W$; $m^{*}_{\pm}$ & densities and scalar densities of the particles of the two halves $W_{1644}=\pm1$ of the field; their differences; the polarization $(n_+-n_-)/n$; the masses of the particles of the two halves & Secs.~\ref{sec:hf}, \ref{sec:kshalves} \\
%%</1644>
$\GS$, $\GV$ & the scalar and the vector coupling of the mean-field description of uniform matter (attraction for $\GS<0$, repulsion for $\GV>0$) & Sec.~\ref{sec:ksmethod} \\
$\hh_{\mathrm{KS}}$; $V_0$; $E^{*}(k)$ & Hamiltonian matrix of a particle in matter; the vector mean field $\GV\,n$; the energy $(k^{2}+\mstq)^{1/2}$ & Eq.~(\ref{eq:hKS}) \\
$D$; $V_D$, $A_D$; $\nu$; $\langle f\rangle$ & number of space dimensions in which the particles have momentum; volume and area of the unit ball in $D$ dimensions; number of internal states of a particle, 8; average over the occupied states & Eq.~(\ref{eq:densities}) \\
$\varepsilon$; $p_j$, $\pint$; $\mu$ & energy density of matter; its pressure along the coordinate $xj$ and the interaction part of the pressure; chemical potential & Eqs.~(\ref{eq:epsKS}), (\ref{eq:pKS}) \\
%%<1671>
$b_t^{\mathrm{ref}}$; $n_{\mathrm{ref}}$ & the scale factor $b_t$ and the density of matter at the reference time $x4_{\mathrm{ref}}$ & Sec.~\ref{sec:ks71} \\
%%</1671>
$\Gamma[S]$, $\Gamma_2$; $S$, $S_0$ & the CJT effective action and its two-particle-irreducible part; full and free propagator of the fermions & Eq.~(\ref{eq:GammaCJT}) \\
$\Omega_{\mathrm{CJT}}$ & the CJT effective potential of uniform matter with the self-interaction & Eq.~(\ref{eq:OmegaCJT}) \\
$\Lambda$; $G_c$ & cutoff of the momenta of the filled states of negative energy (without index or bar; $\Lam$ and $\bar\Lambda$ are cosmological terms); critical coupling of the vacuum & Sec.~\ref{sec:sea} \\
$m_{\mathrm{vac}}$, $m_{\mathrm{bare}}$; $\langle\bar\psi\psi\rangle_{\mathrm{vac}}$ & with the filled states of negative energy: mass of a fermion in vacuum, mass parameter of the Lagrangian, scalar density of the vacuum & Eq.~(\ref{eq:mvac}) \\
$x$; $f_D$ & expectation value of the Vorton-Scalar in units of its vacuum value, $\phi/\vVS$; fraction $m_D/m$ of the fermion mass that is a bare Dirac mass & Eq.~(\ref{eq:unified}) \\
$\vartheta_{\mathrm{VS}}$; $\mth$ & phase of the Vorton-Scalar; mass of its phase mode & Eq.~(\ref{eq:mtheta}) \\
$F_0$, $F_1$, $F_7$, $F_8$; $\Phi$, $\Phi(\vec p)$, $X$ & pairing amplitudes of the four contact channels; a state of a pair of two particles (an antisymmetric $16\times16$ matrix; without suffix, whereas the Vorton-Scalar is $\PhiVS$), the state of the pair along the hidden axis, and its matrix $X=\Cm\gamma^{0}$ & Secs.~\ref{sec:pairs}, \ref{sec:gapcontact} \\
$V_{\mathrm{ct}}$, $V_{\mathrm{VS}}$; $\mct$; $q$ & pairing potentials of the contact interaction and of the exchange of the Vorton-Scalar; range parameter of the contact couplings; momentum transfer & Eqs.~(\ref{eq:Vct}), (\ref{eq:VVS}) \\
$\ma$, $\mb$; $H_{\ma}$, $H_{\mb}$; $H_i$ & the two scale factors of Sec.~\ref{sec:gravity} (sans-serif), for $x0,\dots,x3$ and for $x5,x6,x7$; their rates; the rate of the direction $i$ & Eq.~(\ref{eq:twofactors}) \\
$E^{(k)\mu}{}_{\nu}$; $\alpha_1$, $\alpha_2$, $\alpha_3$; $\bar\Lambda$; $\kappa_8$ & Lovelock tensor of the order $k$ (the order in parentheses; $\EF$, $E_k$, $\Eperp$, $E^{*}$ are energies); couplings of the Lovelock terms (a digit as subscript; $\alpha^{j}$ are matrices); cosmological term; gravitational coupling in eight dimensions & Eq.~(\ref{eq:EL}) \\
$T^{\mu}{}_{\nu}$; loading; $\td^{0}$ & energy-momentum tensor of matter; $\kappa_8\,\varepsilon(0)/(3\,\gff^{2})$; the clock of the undisturbed deflation, $3\,\gff\,x4$ & Sec.~\ref{sec:gravity} \\
\end{tabular}
\end{ruledtabular}
\end{table*}

\begin{table*}
\caption{\label{tab:symbols2}The symbols of this paper, continued: the model and the superfluid.}
\begin{ruledtabular}
\begin{tabular}{p{0.2\textwidth}p{0.6\textwidth}p{0.14\textwidth}}
symbol & meaning & defined in \\
\hline
$\mathcal L_{\tg}^{\mathrm{VS}}$; $\PhiVS$, $\hat\Phi_{\mathrm{VS}}$ & Lagrangian of the model; the Vorton-Scalar and its weighted form $\sqrt{\Vd}\,\PhiVS$ & Eq.~(\ref{eq:LVS}) \\
$\lamVS$, $\vVS$, $\yVS$ & self-coupling of the Vorton-Scalar, the constant that fixes its vacuum value, its Yukawa coupling & Sec.~\ref{sec:clock} \\
$U(1)_F$, $\UVS$; $\epsilon$ & the symmetries of fermion number and of chirality; parameter of a symmetry transformation & Sec.~\ref{sec:clock} \\
$m$, $m_D$, $m_Y$; $m_s$; $\gVS$, $\aVS$ & mass of a fermion in vacuum, $m=m_D+m_Y$: the bare Dirac mass of the Lagrangian plus the part $m_Y=\yVS\vVS/\sqrt2$ from the Vorton-Scalar; mass of the radial mode; coupling $\gVS=m_Y/\vVS$ of both modes to the fermions, and $\gVS^{2}/4\pi$ & Eq.~(\ref{eq:masses}) \\
$\DVS$, $C$; $D_c$ & the two numbers that decide the state of matter; critical value of $\DVS$ for a bound pair & Eq.~(\ref{eq:DC}) \\
$\td$; $\tpair$ & the deflation clock; the clock measured from the binding of the pair & Eqs.~(\ref{eq:clock}), (\ref{eq:tpair}) \\
$\varphi_{\mathrm{bos}}$, $m_{\mathrm{bos}}$, $\varpi_b$ & a mode of the Vorton-Scalar, its mass, the frequency of its weighted form & Sec.~\ref{sec:scalarmodes} \\
$F_S$, $F_P$; $W(k,k')$ & exchange amplitudes of the radial and of the phase mode; kernel of the gap equation (always written with its two arguments) & Sec.~\ref{sec:gap} \\
$\lambda$, $\lambda_S$, $\lambda_P$; $\Delta_F$; $T_c$ & pairing strength, and its parts from the radial and from the phase mode; gap at the Fermi momentum; critical temperature & Sec.~\ref{sec:gap} \\
$\kF$, $\EF$, $v_F$, $N(0)$ & Fermi momentum, Fermi kinetic energy, Fermi velocity, density of states at the Fermi surface per internal state & Secs.~\ref{sec:gap}, \ref{sec:cjt} \\
$a_{\mathrm{sc}}$, $E_b$ & scattering length of two fermions; binding energy of the pair & Sec.~\ref{sec:pair} \\
$\Xid$; $\Deq$; $\Esh$; $\nex$ & the parameter that decides whether the paired state follows the deflation; equilibrium gap; half-width of the shell of states; density of broken pairs & Sec.~\ref{sec:quench} \\
$\xi_k$, $E_k$, $\kappa$, $n_k$ & single-particle energy, quasi-particle energy, pair correlation and occupation of a pair of states & Eq.~(\ref{eq:bloch}) \\
$V[\Delta]$; $\phi$; $\mst$; $T_{1/2}$ & effective potential of the paired state; expectation value of the Vorton-Scalar, $\langle\PhiVS\rangle=\phi/\sqrt2$; mass of a fermion in matter; temperature at which the fermions have lost half of their mass & Secs.~\ref{sec:cjt}, \ref{sec:ks} \\
$a(T)$, $b$, $K_{\mathrm{GL}}$; $T_{\mathrm{high}}$, $T_0$ & coefficients of the Ginzburg-Landau energy of the paired state near $T_c$ (quadratic, quartic and gradient term); a closed estimate of the temperature above which the expectation value of the Vorton-Scalar vanishes, and the temperature at which it has fallen to 0.02 of its vacuum value & Appendix~\ref{app:tables} \\
$\Delta^{a}$, $\Delta_t$; $\MM$; $d_\pm$ & order parameter (frame components), and its coordinate components along the deflating coordinates; gap matrix; the two fourfold gaps of a state & Sec.~\ref{sec:order} \\
$SU(2)_A$, $SU(2)_B$, $U(1)_E$\SIG{; $SO(3)_T$}{} & the factors of the internal symmetry of the fields that do not depend on the extra coordinates\SIG{; the rotations of the three additional times $x5,x6,x7$}{} & Sec.~\ref{sec:hidden} \\
$P$, $P_c$; $w_z$; $\theta_{12}$ & generator of $U(1)_E$ and the product of the five matrices of the sector; spinor of a zero mode; angle around a vortex line & Sec.~\ref{sec:vortex} \\
\end{tabular}
\end{ruledtabular}
\end{table*}

\begin{table*}
\caption{\label{tab:symbols3}The symbols of this paper, continued: the cosmological part and the
Higgs hypothesis.}
\begin{ruledtabular}
\begin{tabular}{p{0.2\textwidth}p{0.6\textwidth}p{0.14\textwidth}}
symbol & meaning & defined in \\
\hline
$a_{\mathrm{sp}}$, $H$, $H_0$; $\GN$, $G_8$ & scale factor of ordinary space, Hubble rate, Hubble constant; Newton's constant and the constant of gravitation in eight dimensions & Sec.~\ref{sec:cosmo} \\
$\rho$, $\rho_m$; $\omega_b$, $\omega_c$ & energy density in ordinary space, and that of matter; the densities of baryons and of cold dark matter of the standard model of cosmology & Sec.~\ref{sec:cosmo} \\
$\theta_*$, $r_s$; $t_0$; $z_*$, $\YHe$ & angle under which the sound horizon at last scattering is seen, and that sound horizon; age of the universe; redshift of last scattering, mass fraction of helium & Sec.~\ref{sec:cosmo} \\
$\delta_m$, $f_g$, $\sigma_8$ & density contrast of matter, its growth rate, amplitude of clustering & Sec.~\ref{sec:cosmo} \\
$p_{\mathrm{agent}}$ & pressure, in ordinary space, of the agent: of whatever keeps the three rates of deflation constant & Eq.~(\ref{eq:pagent}) \\
$\tcf$; $\Hcf$, $\He$ & conformal time of ordinary space (a prime in Sec.~\ref{sec:pert} is the derivative with respect to it); conformal Hubble rate, and $\Hcf-a_{\mathrm{sp}}\gsum/2$ & Sec.~\ref{sec:pert} \\
$k$ (Sec.~\ref{sec:cosmo}); $\ell$, $D_\ell^{TT}$ & comoving wavenumber of a perturbation of ordinary space; multipole, and the temperature spectrum of the microwave background & Secs.~\ref{sec:pert}, \ref{sec:spectra} \\
$\hsy$, $\etasy$; $\alsy$ & the two functions of the perturbed metric of ordinary space in the synchronous gauge; $(\hsy'+6\etasy')/2k^{2}$ & Sec.~\ref{sec:pert} \\
$\cD$, $\cT$, $\cP$, $\cA$ & $8\pi\GN a_{\mathrm{sp}}^{2}$ times the perturbations of the energy density, of the divergence of the momentum density, of the pressure, and the anisotropic stress, of matter and radiation together & Eq.~(\ref{eq:pert}) \\
$\cE$, $\cE_P$, $\cE_m$, $\cE_r$; $\cB$, $\cZ$ & $8\pi\GN a_{\mathrm{sp}}^{2}$ times the density and the pressure of the unperturbed universe, and the parts of the first from matter and from radiation; two combinations that vanish without deflation & Eq.~(\ref{eq:Z}) \\
$\cYM$, $\cYT$, $\cYS$ & the gaps: what the momentum constraint, the trace equation and the equation of the shear lack with matter and radiation alone & Eq.~(\ref{eq:gaps}) \\
$\mathcal{U}$; $R_{\mathrm{TF}}$ & $\etasy+\alsy'$, the sum of the two gravitational potentials of the conformal Newtonian gauge; the radius of a core of condensed pairs that is supported by its own pressure & Appendix~\ref{app:tables} \\
$\Delta\chi^{2}$; $N$, $z_t$; $c_{s0}$ & difference of $\chi^{2}$ from the best fit of the standard model to the mock data; density of extra radiation in units of that of one neutrino species, and the transition redshift of a stepped fluid; sound speed today of a condensate of bound pairs & Secs.~\ref{sec:forecast}, \ref{sec:dm} \\
$\sigma_v$; $m_{\mathrm{deg}}$, $m_{B,\max}$, $m_B$; $\kappa_{\mathrm{SI}}$ & velocity dispersion of the dark matter of a halo; the largest masses at which fermions are degenerate and bosons of the mass $m_B$ condensed there; the bound on the cross-section per mass of dark matter & Sec.~\ref{sec:dm} \\
$\vEW$, $\nEW$; $\alpha_c$, $\varrho_c$ & electroweak scale and the exponent with which it would follow the coupling under H2; critical coupling of the vacuum condensate, and an abbreviation & Sec.~\ref{sec:higgs} \\
$m_p/m_e$; $G_{\mathrm{Fermi}}$; $\alpha_{\mathrm{em}}$ & ratio of the proton mass to the electron mass; Fermi constant; fine-structure constant & Sec.~\ref{sec:cosmo} \\
H1, H2, H3; A1, A2, A3; HV & the three hypotheses about the superfluid; the three assumptions of the cosmological part (A3, on the perturbations, with its closures ``Momentum'', ``Pressure'' and ``Shear''); the assumption that the vacuum value of the Vorton-Scalar sets the electroweak scale & Secs.~\ref{sec:intro}, \ref{sec:cosmo}, \ref{sec:pert} \\
\end{tabular}
\end{ruledtabular}
\end{table*}

\section{The Dirac matrices}
\label{app:matrices}

Table~\ref{tab:gammas} prints the eight Dirac matrices of the orthonormal frame and three matrices
built from them. All are real, with entries $0$, $+1$ and $-1$, and each Dirac matrix is a signed
permutation matrix.
%%<1644>
The matrices $\gamma^{0},\dots,\gamma^{3}$ of the space-like coordinates are symmetric and square
to $+1$; the matrices $\gamma^{4},\dots,\gamma^{7}$ of the time-like coordinates are antisymmetric
and square to $-1$ [Eq.~(\ref{eq:clifford})]. Their product
$\omega_{1644}=\gamma^{0}\gamma^{1}\cdots\gamma^{7}$ is the chirality matrix $\chi$; it is real and
diagonal, $\mathrm{diag}(-1_8,+1_8)$, so that the chiral projectors $P^{L,R}=\tfrac12(1\mp\chi)$
select the upper and the lower eight components. The matrix $\Cm=\gamma^{0}\gamma^{1}\gamma^{2}\gamma^{3}$
is real and symmetric and turns each Dirac matrix into minus its transpose,
$\Cm\gamma^{A}(\Cm)^{-1}=-(\gamma^{A})^{T}$; it defines the covariant adjoint
$\tilde\psi=\psi^\dagger\Cm$. The matrix $C^{+}=\Cm\omega_{1644}$ turns each Dirac matrix into plus
its transpose. Three further matrices follow by one multiplication: the Hermitian matrices
$\beta=-i\gamma^{4}$ and $\alpha^{j}=-\gamma^{4}\gamma^{j}$ ($j=0,1,2,3$) of the quantized sector,
and $W_{1644}=i\,\Cm\gamma^{4}$, which is Hermitian, squares to one, commutes with $\beta$ and with
the $\alpha^{j}$, and satisfies $\Cm=-W_{1644}\beta$.
%%</1644>
%%<1671>
The matrix $\gamma^{4}$ of the time is symmetric and squares to $+1$; the seven matrices of the
space-like coordinates are antisymmetric and square to $-1$ [Eq.~(\ref{eq:clifford})]. Their
product $\omega_{1671}=\gamma^{0}\gamma^{1}\cdots\gamma^{7}$ is real and squares to $-1$; the
chirality matrix is $\chi=i\,\omega_{1671}$, and the chiral projectors $P^{L,R}=\tfrac12(1\mp\chi)$
are complex. The matrix $\Cm=\gamma^{4}\omega_{1671}$ turns each Dirac matrix into minus its
transpose, $\Cm\gamma^{A}(\Cm)^{-1}=-(\gamma^{A})^{T}$, and enters the pair operators
$\psi^{T}\Cm\gamma^{A}\psi$; the matrix $C^{+}=\gamma^{4}$ turns each into plus its transpose. Two
kinds of matrices follow by one multiplication: $\beta=\gamma^{4}$, which defines the adjoint
$\bar\psi=\psi^\dagger\beta$, and the real symmetric matrices $\alpha^{A}=\gamma^{4}\gamma^{A}$
($A\ne4$) of the Hamiltonian.
%%</1671>
The program verifies every property stated here in exact integer arithmetic
(Appendix~\ref{app:checks}).

\begin{table*}
\caption{\label{tab:gammas}The real matrices of the \SIG{$4+4$}{$7+1$} spacetime ($+$ stands for
$+1$, $-$ for $-1$, a dot for $0$; the lines separate the upper and the lower eight components).
First and second row: the eight Dirac matrices $\gamma^{A}$ of the orthonormal frame. Third row:
their product $\omega_{\SIG{1644}{1671}}$, and the matrices $\Cm$ and $C^{+}$.}
\setlength{\tabcolsep}{1pt}
\begin{ruledtabular}
\begin{tabular}{cccc}
$\gamma^{0}$ & $\gamma^{1}$ & $\gamma^{2}$ & $\gamma^{3}$ \\[2pt]
{\fontsize{6}{7.2}\selectfont$\displaystyle{{MATRIX:gamma\SIG{1644}{1671}[0]}}$} &
{\fontsize{6}{7.2}\selectfont$\displaystyle{{MATRIX:gamma\SIG{1644}{1671}[1]}}$} &
{\fontsize{6}{7.2}\selectfont$\displaystyle{{MATRIX:gamma\SIG{1644}{1671}[2]}}$} &
{\fontsize{6}{7.2}\selectfont$\displaystyle{{MATRIX:gamma\SIG{1644}{1671}[3]}}$} \\[14pt]
$\gamma^{4}$ & $\gamma^{5}$ & $\gamma^{6}$ & $\gamma^{7}$ \\[2pt]
{\fontsize{6}{7.2}\selectfont$\displaystyle{{MATRIX:gamma\SIG{1644}{1671}[4]}}$} &
{\fontsize{6}{7.2}\selectfont$\displaystyle{{MATRIX:gamma\SIG{1644}{1671}[5]}}$} &
{\fontsize{6}{7.2}\selectfont$\displaystyle{{MATRIX:gamma\SIG{1644}{1671}[6]}}$} &
{\fontsize{6}{7.2}\selectfont$\displaystyle{{MATRIX:gamma\SIG{1644}{1671}[7]}}$} \\[14pt]
$\omega_{\SIG{1644}{1671}}$ & $\Cm$ & $C^{+}$ & \\[2pt]
{\fontsize{6}{7.2}\selectfont$\displaystyle{{MATRIX:omega\SIG{1644}{1671}}}$} &
{\fontsize{6}{7.2}\selectfont$\displaystyle{{MATRIX:Cm\SIG{1644}{1671}}}$} &
{\fontsize{6}{7.2}\selectfont$\displaystyle{{MATRIX:Cp\SIG{1644}{1671}}}$} & \\
\end{tabular}
\end{ruledtabular}
\end{table*}

\section{Curvature of the deflating metric}
\label{app:curvature}

This appendix derives the Einstein tensor and the Lanczos tensor of Sec.~\ref{sec:geometry} in a
form that can be checked by hand, and gives the coordinates and the additional symmetries of the
case of equal rates.

\subsection{Curvatures on pairs of directions}

In the orthonormal frame the Riemann tensor (\ref{eq:riemann}) of the deflating metric is diagonal
on pairs of directions: for each pair of frame directions $A\ne B$ it has one independent component,
$K_{AB}=K_{BA}=R^{AB}{}_{AB}$ (no summation), the sectional curvature of the plane that the two
directions span. With the Riemann tensor (\ref{eq:riemann}) and the flat metric (\ref{eq:etam}) the
sectional curvatures are, on the $4+4$ and on the $7+1$ spacetime alike,
\begin{equation}
K_{4t}=\rt^{2},\qquad K_{ts}=\rt\,\mathrm{rate}_s\quad(t\ne s),
\label{eq:sectional}
\end{equation}
and $K_{AB}=0$ for every pair that contains one of the flat directions $0,1,2,3$. \SIG{(A sectional
curvature $K_{AB}$ carries two subscripts; it is not the matrix $K$ of Sec.~\ref{sec:dirac}.) }{}For a
Riemann tensor of this kind the mixed components of the Einstein tensor and of the Lanczos tensor
are diagonal, and
\begin{align}
\cG^{A}{}_{A}&=-\sum_{\{B,C\}}K_{BC},\label{eq:pairG}\\
\cH^{A}{}_{A}&=-4\sum_{\{B,C\},\{B',C'\}}K_{BC}\,K_{B'C'},
\label{eq:pairH}
\end{align}
without summation over $A$. The first sum runs over the pairs of directions that do not contain the
direction $A$. The second runs over the unordered couples of two such pairs that have no direction
in common; a set of four directions contributes its three divisions into two pairs. As a test, for a
space of constant curvature in eight dimensions, where all $K_{BC}$ equal one number, the first sum
has $21$ terms and the second $3\times35$, so that the two formulas give $-21$ times that number and
$-420$ times its square, the known values $-\tfrac12\cdot7\cdot6$ and
$-\tfrac12\cdot7\cdot6\cdot5\cdot4$.

\emph{The Einstein tensor.} With Eq.~(\ref{eq:sectional}), Eq.~(\ref{eq:pairG}) gives
\begin{equation}
\begin{aligned}
\cG^{4}{}_{4}&=-\gprod,\qquad\cG^{j}{}_{j}=-(\gsq+\gprod),\\
\cG^{5}{}_{5}&=-(\gss^{2}+\gvv^{2}+\gss\,\gvv),
\end{aligned}
\end{equation}
and the corresponding expressions for $\cG^{6}{}_{6}$ and $\cG^{7}{}_{7}$. Because
$\gsum^{2}=\gsq+2\,\gprod$, these are Eq.~(\ref{eq:einstein}).

\emph{The Lanczos tensor.} A couple of pairs contributes to Eq.~(\ref{eq:pairH}) only if neither
pair contains a flat direction, that is, only if the four directions are $4,5,6,7$. For a flat
direction $A=j$ these are allowed, and their three divisions into two pairs give
\begin{equation}
K_{45}K_{67}+K_{46}K_{57}+K_{47}K_{56}=\gff\,\gss\,\gvv\,\gsum ,
\end{equation}
which is Eq.~(\ref{eq:lanczos}). For $A=4,5,6,7$ the direction $A$ is itself among the four, and
these components vanish. The Lovelock tensor of the third order is built in the same way from sets
of six directions, and every such set contains a flat direction: it vanishes.

\emph{The field equations.} In Eq.~(\ref{eq:egb}) the components along $x4$ and along the deflating
coordinates read $\cG^{4}{}_{4}+\Lam=0$ and $\cG^{t}{}_{t}+\Lam=0$. Since
$\cG^{5}{}_{5}-\cG^{6}{}_{6}=(\gff-\gss)\,\gsum$, they can hold together only if the three rates are
equal, and then $\Lam=3\,\gff^{2}$. The components along the flat coordinates read
$-6\,\gff^{2}+\Lam-12\,\aGB\,\gff^{4}=0$, which gives $\aGB=-1/(4\,\gff^{2})$:
Eq.~(\ref{eq:egbsol}). The program verifies Eqs.~(\ref{eq:pairG}) and (\ref{eq:pairH}) for the
deflating metric in exact arithmetic, and an independent program obtains the same Lanczos tensor
from its standard definition by computer algebra, for arbitrary rates and constants.

\subsection{Equal rates: a space of constant curvature}

Let $\gff=\gss=\gvv$, and introduce in place of $x4,\dots,x7$ the coordinates
\begin{equation}
\bar x4=\frac{e^{\gff\,x4}}{\gff},\qquad\bar xt=\ct\,xt\quad(t=5,6,7).
\end{equation}
The part of the line element that contains $x4,\dots,x7$ becomes
\begin{equation}
-dx4^{2}\SIG{-}{+}e^{-2\,\gff\,x4}\sum_t d\bar xt^{2}
=\frac{-d\bar x4^{2}\SIG{-}{+}\sum_t d\bar xt^{2}}{(\gff\,\bar x4)^{2}} .
\label{eq:conformal}
\end{equation}
\SIG{This is minus the line element of four-dimensional hyperbolic space of the
radius $1/\gff$, in its upper-half-space form: the four time-like coordinates span a
negative-definite space of constant curvature.}{This is the line element of de
Sitter space of the radius $1/\gff$, in the half of its flat slicing in which space contracts. De
Sitter space is familiar as the geometry of inflation, where the same line element is used with the
opposite direction of time.} The line element (\ref{eq:conformal}) is a multiple of a flat one. The
special conformal transformations of that flat metric along the three axes $\bar x5,\bar x6,\bar x7$
leave the factor $1/\bar x4^{2}$ unchanged and are therefore symmetries. Their generators are, for
$s=5,6,7$,
\begin{equation}
\begin{aligned}
\cX_{(s)}={}&2\,\bar xs\,\Big[\bar x4\,\frac{\partial}{\partial\bar x4}+\sum_t\bar xt\,\frac{\partial}{\partial\bar xt}\Big]\\
&-\Big(\sum_t\bar xt^{2}\SIG{+}{-}\bar x4^{2}\Big)\frac{\partial}{\partial\bar xs},
\end{aligned}
\label{eq:special}
\end{equation}
with $\bar x4\,\partial/\partial\bar x4=\gff^{-1}\partial_4$. Together with the 14 Killing vectors of
the general case and the three rotations of $x5,x6,x7$ they are the 20 Killing vectors of the case
of equal rates.

